\begin{abstract}
La creciente popularidad de las nubes de puntos 3D en aplicaciones como la realidad virtual y los vehículos autónomos
han generado una necesidad crítica de algoritmos de compresión eficientes. Este trabajo explora un método de compresión
con pérdida para nubes de puntos estáticas, utilizando la Transformada de Fourier sobre Grafos (GFT) para explotar la
correlación espacial de los atributos de la nube, como la luminancia. El rendimiento de la GFT depende
fundamentalmente de la estructura del grafo subyacente. Para optimizar esta dependencia, se propone un novedoso
algoritmo de agrupamiento basado en optimización Tasa-Distorsión (RD). A diferencia del k-means tradicional, este
método agrupa bloques de la nube de puntos asignándolos a centroides (modos de la señal) que minimizan un costo
combinado de tasa de bits y error de cuantización. Se introduce un esquema de Simulated Annealing para mejorar la
exploración del espacio de soluciones y evitar la convergencia a mínimos locales, resultando en un conjunto de
transformadas adaptativas que mejoran la eficiencia de la compresión.
\end{abstract}

\begin{IEEEkeywords}
Nubes de Puntos, Compresión de Datos, Transformada de Fourier sobre Grafos, Optimización Tasa-Distorsión, Simulated
Annealing.
\end{IEEEkeywords}

\section{Introducción}
Las nubes de puntos 3D se han convertido en un formato de datos fundamental para una amplia gama de tecnologías
emergentes. Sin embargo, su representación no estructurada y de alta densidad resulta en archivos de gran tamaño, lo
que presenta desafíos significativos para su almacenamiento y transmisión. La compresión de nubes de puntos es, por lo
tanto, un área de investigación activa y crucial.

Los métodos de compresión basados en transformadas, como la GFT, han demostrado ser prometedores al adaptar la
transformada a la geometría irregular de la nube de puntos. La GFT proyecta la señal (atributos de los puntos) sobre
una base de vectores propios de la matriz Laplaciana del grafo, logrando una representación compacta en el dominio
espectral. La eficiencia de esta compactación está intrínsecamente ligada a la definición del grafo. Un grafo bien
construido, que capture las correlaciones de la señal, resultará en coeficientes GFT más esparcidos y, por lo tanto,
más compresibles.

Este proyecto aborda el problema de encontrar un conjunto óptimo de grafos (y sus correspondientes GFTs) mediante un
método de agrupamiento que no solo considera la similitud de la señal, sino también el costo final de compresión.

\section{Preliminares o Marco Teórico}
\begin{definition}[Procesamiento de Señales sobre Grafos]
El Procesamiento de Señales sobre Grafos (GSP) extiende los conceptos del procesamiento de señales tradicional a datos
definidos en dominios irregulares. Una señal sobre un grafo $\mathcal{G} = (\mathcal{V}, \mathcal{E})$ es una función
$f: \mathcal{V} \to \mathbb{R}$ que asigna un valor a cada vértice. La estructura del grafo se representa mediante su
matriz Laplaciana, $L$.
\end{definition}

\begin{definition}[Transformada de Fourier sobre Grafos (GFT)]
La GFT de una señal $f$ se define como su proyección sobre los vectores propios de la matriz Laplaciana del grafo. Si
$L = U \Lambda U^T$ es la descomposición propia de $L$, la GFT de $f$ es $\hat{f} = U^T f$. La GFT inversa es $f = U
\hat{f}$. La GFT concentra la energía de la señal en unos pocos coeficientes de "baja frecuencia" si la señal es suave
con respecto a la estructura del grafo.
\end{definition}

\begin{definition}[Optimización Tasa-Distorsión (RD)]
En compresión con pérdida, existe un compromiso entre la tasa de bits (Rate, R) y la distorsión (Distortion, D). La
optimización Tasa-Distorsión busca minimizar la distorsión para una tasa de bits dada, o viceversa. Formalmente, se
busca minimizar el costo Lagrangiano $J = D + \lambda R$, donde $\lambda$ es un multiplicador de Lagrange que controla
el compromiso entre tasa y distorsión.
\end{definition}

\section{Metodología o Configuración Experimental}
La solución propuesta se basa en un algoritmo de agrupamiento iterativo que optimiza la elección de la GFT para cada
bloque de la nube de puntos.

\subsection{Particionamiento y Aproximación de Señal}
Inicialmente, la nube de puntos se divide en bloques cúbicos utilizando una curva de Morton para preservar la
localidad espacial. Para cada bloque, se aproxima el atributo de luminancia de los puntos mediante un modelo lineal.
La pendiente de este modelo (un vector 3D) sirve como un descriptor compacto de la tendencia de la señal dentro del
bloque.

\subsection{Agrupamiento Tasa-Distorsión (RD Clustering)}
El corazón de nuestra metodología es un algoritmo de agrupamiento que refina un conjunto de centroides (pendientes) de
forma iterativa. A diferencia de k-means, el objetivo es minimizar el costo RD global.

\begin{enumerate}
\item \textbf{Inicialización}: Se inicializa un conjunto de $k$ centroides. Uno de ellos es fijo y representa una
pendiente nula (transformada estructural), mientras que los otros se inicializan a partir de las pendientes de los
bloques.

\item \textbf{Paso de Asignación Estocástica}: En cada iteración, para cada bloque, se calcula el costo RD que
resultaría de asignarlo a cada uno de los $k$ clusters. Cada cluster $j$ tiene asociado un centroide (pendiente)
$\vec{s}_j$, que a su vez define una GFT adaptativa $\mathcal{F}_j$. El costo se calcula como $J_j = D_j + \lambda
R_j$, donde $D_j$ es el error de cuantización y $R_j$ es una estimación de la tasa de bits de los coeficientes
$\mathcal{F}_j(f)$.

Para evitar la convergencia a mínimos locales, la asignación no es determinista. Se utiliza un esquema de Simulated
Annealing, donde la probabilidad de asignar un bloque al cluster $j$ es proporcional a $e^{-J_j/T}$. La temperatura
$T$ disminuye gradualmente en cada iteración, permitiendo la exploración al principio y la explotación al final.

\item \textbf{Paso de Actualización de Centroides}: Una vez que todos los bloques han sido reasignados, se recalcula
el centroide de cada cluster como la pendiente promedio de todos los bloques que le pertenecen.
\end{enumerate}

\begin{algorithm}[H]
\caption{Algoritmo de Agrupamiento Tasa-Distorsión (RD Clustering)}
\begin{algorithmic}[1]
\REQUIRE Nube de puntos $P$, número de clusters $K$, parámetros de Simulated Annealing ($T_{inicial}, T_{final}, \alpha$), $\lambda$ schedule, $self\_loop\_threshold$, $self\_loop\_weight$.
\ENSURE Asignaciones de clusters para cada bloque $L$, centroides de clusters $S$.

\STATE Inicializar $S_0 \leftarrow \text{pendiente nula}$
\STATE Inicializar $S_1, \dots, S_{K-1} \leftarrow \text{pendientes aleatorias o K-means inicial}$
\STATE Inicializar $L \leftarrow \text{asignaciones aleatorias de bloques a clusters}$
\STATE $T \leftarrow T_{inicial}$
\STATE Precomputar GFT estructural para todos los bloques y almacenar en caché

\FOR{iteración = 1 \TO max\_iteraciones}
    \FOR{cada bloque $b$ en $P$}
        \STATE $costos \leftarrow []$
        \FOR{cada cluster $j = 0 \TO K-1$}
            \IF{$j = 0$}
                \STATE $coeffs_j \leftarrow \text{GFT estructural de } b$
            \ELSE
                \STATE $graph_j \leftarrow \text{construir AttributeGraph para } b \text{ con } S_j, self\_loop\_threshold, self\_loop\_weight$
                \STATE $coeffs_j \leftarrow \text{calcular GFT adaptativa de } b \text{ usando } graph_j$
            \ENDIF
            \STATE $costos.append(\text{calcular costo RD}(coeffs_j, \lambda))$
        \ENDFOR
        \STATE $L_b \leftarrow \text{seleccionar cluster de } costos \text{ usando Simulated Annealing con } T$
    \ENDFOR
    \STATE Actualizar $S_j \leftarrow \text{pendiente promedio de bloques asignados a cluster } j$
    \STATE $T \leftarrow T \cdot \alpha$ (enfriamiento)
    \STATE Actualizar $\lambda$ si es necesario
    \IF{criterio de convergencia se cumple}
        \STATE \textbf{BREAK}
    \ENDIF
\ENFOR
\RETURN $L, S$
\end{algorithmic}
\end{algorithm}

\subsection{GFT Adaptativa}
Cada centroide define una GFT. El centroide nulo corresponde a una GFT estructural, basada únicamente en la distancia
euclidiana. Los centroides dinámicos definen un "AttributeGraph", donde la estructura del grafo se modifica añadiendo
auto-bucles (self-loops) a los nodos que actúan como "sumideros" en el campo de gradiente definido por la pendiente
del centroide. Esto adapta la transformada a la tendencia de la señal, mejorando la compactación de energía.

\subsection{Parámetros del Experimento}
El experimento se realizó sobre la nube de puntos "longdress" del dataset 8iVFB v2, utilizando el frame 1051. Los
parámetros clave del algoritmo se configuraron de la siguiente manera:
\begin{itemize}
    \item \textbf{Número de clusters (k)}: 4 (1 estructural y 3 dinámicos).
    \item \textbf{Self-loop threshold}: 0.5.
    \item \textbf{Self-loop weight}: 1.6.
    \item \textbf{Learning rate (actualización de centroides)}: 0.5.
    \item \textbf{Parámetros de Simulated Annealing}:
    \begin{itemize}
        \item Temperatura inicial: 1.0.
        \item Temperatura final: 0.1.
        \item Tasa de enfriamiento: 0.95.
    \end{itemize}
\end{itemize}

\section{Resultados y Análisis}
El algoritmo se ejecutó durante 17 iteraciones. A continuación se presenta un análisis de los resultados obtenidos a partir de las métricas y visualizaciones generadas.

\subsection{Utilización de Clusters}
La Figura \ref{fig:blocks_per_cluster} muestra la evolución del número de bloques asignados a cada cluster. Se observa que los cuatro clusters fueron utilizados durante la ejecución. Sin embargo, el cluster 0 (estructural) y el cluster 1 (dinámico) son claramente dominantes, agrupando la gran mayoría de los bloques. Los clusters 2 y 3, aunque reciben asignaciones, mantienen una población significativamente menor a lo largo de las iteraciones. Esto sugiere que solo un modo dinámico (Cluster 1) es fuertemente preferido por el algoritmo.

\begin{figure}[H]
    \centering
    % Placeholder for blocks_per_cluster_vs_iteration.png
    \includegraphics[width=0.8\textwidth]{path/to/your/images/blocks_per_cluster_vs_iteration.png}
    \caption{Evolución del número de bloques por cluster.}
    \label{fig:blocks_per_cluster}
\end{figure}

\subsection{Análisis de la Ganancia por Cluster}
La Figura \ref{fig:gain_per_cluster} presenta la ganancia media de cada cluster dinámico. Esta métrica es crucial, ya que indica el beneficio en costo RD de usar una GFT adaptativa en lugar de la estructural. El análisis revela un resultado clave:
\begin{itemize}
    \item El \textbf{Cluster 1} muestra una ganancia media consistentemente positiva y creciente, validando la efectividad de su GFT adaptativa para un subconjunto de bloques.
    \item Los \textbf{Clusters 2 y 3} tienen una ganancia media de cero en todas las iteraciones. Esto significa que, a pesar de que se les asignan bloques (debido a la exploración estocástica del Simulated Annealing), sus GFTs adaptativas no logran reducir el costo RD en comparación con la GFT estructural.
\end{itemize}
Este hallazgo es el más crítico del experimento, ya que demuestra que el algoritmo, bajo la configuración actual, no es capaz de encontrar más de un modo dinámico que sea beneficioso para la compresión.

\begin{figure}[H]
    \centering
    % Placeholder for gain_per_cluster_vs_iteration.png
    \includegraphics[width=0.8\textwidth]{path/to/your/images/gain_per_cluster_vs_iteration.png}
    \caption{Evolución de la ganancia media por cluster dinámico.}
    \label{fig:gain_per_cluster}
\end{figure}

\subsection{Costo RD y Entropía}
La Figura \ref{fig:total_cost} muestra el costo RD total a lo largo de las iteraciones. El comportamiento escalonado creciente es el esperado, ya que el parámetro $\lambda$ del costo Lagrangiano se incrementa, forzando al algoritmo a priorizar una menor tasa de bits (más compresión) a expensas de una mayor distorsión. La Figura \ref{fig:entropy} muestra la entropía de la distribución de clusters. La tendencia general a la baja indica que las asignaciones de los bloques a los clusters se vuelven más definidas y menos aleatorias a medida que el algoritmo converge y la temperatura del Simulated Annealing disminuye.

\begin{figure}[H]
    \centering
    % Placeholder for total_cost_vs_iteration.png
    \includegraphics[width=0.8\textwidth]{path/to/your/images/total_cost_vs_iteration.png}
    \caption{Evolución del costo RD total.}
    \label{fig:total_cost}
\end{figure}

\begin{figure}[H]
    \centering
    % Placeholder for entropy_vs_iteration.png
    \includegraphics[width=0.8\textwidth]{path/to/your/images/entropy_vs_iteration.png}
    \caption{Evolución de la entropía de los clusters.}
    \label{fig:entropy}
\end{figure}

\section{Conclusiones y Trabajo Futuro}
El algoritmo de RD Clustering propuesto demuestra ser capaz de identificar un modo de señal dinámico (Cluster 1) que
mejora significativamente la eficiencia de la compresión en comparación con una GFT puramente estructural.
La
metodología de GFT adaptativa, basada en la adición de auto-bucles, es por tanto validada como una técnica prometedora.

Sin embargo, el experimento también revela una limitación importante: la dificultad del algoritmo para encontrar
múltiples clusters dinámicos beneficiosos. La nula ganancia de los clusters 2 y 3 sugiere que la inicialización de
sus centroides, o la propia capacidad del modelo de pendiente lineal, no son suficientes para capturar otros modos de
señal relevantes en el dataset.

Como trabajo futuro, se proponen las siguientes líneas de investigación:
\begin{enumerate}
    \item \textbf{Reducción del número de clusters}: Realizar una ejecución con un menor número de clusters (e.g., 2)
para forzar al algoritmo a encontrar el único cluster dinámico más representativo.
    \item \textbf{Optimización de hiperparámetros}: Realizar una búsqueda sistemática de los parámetros
`self_loop_threshold` y `self_loop_weight` para mejorar la efectividad de la GFT adaptativa.
    \item \textbf{Modelos de señal más complejos}: Investigar el uso de modelos de mayor grado (e.g., polinomios
cuadráticos) para la aproximación de la luminancia, lo que podría permitir capturar variaciones más complejas y,
por tanto, descubrir más clusters dinámicos útiles.
\end{enumerate}

\end{document}